%% kapa: exam=1 round=22 year=2011 period=1 form=A subject=economics subject_ko=경제원론 seq=6 number=46
%% source: Q-Net 감정평가사 기출문제 (cst00309) artlSeq=5036790 file=extracted/1008/5036790/2011년 제22회 감정평가사 1차시험 기출문제/1교시.hwp
%% conversion: kordoc 4.24.51 -> tools/gen1.py (auto); review=proofread-sonnet; flags=none
\begin{question}{46}
소비자 선호체계와 소비자 선택에 관한 설명으로 옳지 않은 것은?
\begin{choices}
  \item 효용함수가 U=X+Y이고, X재의 가격이 Y재의 가격보다 높을 때 X재만을 소비한다.
  \item 효용함수가 U=min\{X, Y\} 이라면 항상 동일한 양의 X재와 Y재를 소비한다.
  \item 한계대체율은 무차별곡선 기울기의 절대값을 나타낸다.
  \item 두 무차별곡선이 교차할 수 없다는 성질은 선호체계의 이행성으로부터 도출된다.
  \item 효용함수가 U=(X+Y)\textsuperscript{2}이면, 무차별곡선은 직선이다.
\end{choices}
\end{question}
\answer{1}
